\chapter{Introduction}
\label{sec:1}
\begin{verse}
I used to want to get \\
the chance to show \\
the world I'm smart (ha)

Isn't that dumb? \\
I should have focused \\
mostly on the heart

'Cause i've seen smarter \\
people trample life \\
like it's an art
\end{verse}
\begin{flushright}\small\itshape
a few words for the firing squad (radiation) \\
Run the Jewels 4
\end{flushright}
“Intellect estranged from compassion” characterizes many of the twentieth century's worst actors. This phrase also describes what our industries are building today: machines that can do almost anything, but are indifferent to whether those deeds help or harm humanity. This book asks what it would take to build a machine of that capability that is altruistic and antifascist by design — one that stands, by definition, for humanity and against concentrated, unaccountable power.

The word that organizes this field is safety, and it is doing less work than it appears to. A safe system is safe \emph{for} someone: it does what we say, and it stops when whoever holds it says stop. That property is indexed to the party in possession, by construction. Ethics is indexed too, and to somebody else — the party who can be wronged, who is rarely the party holding the switch. A system that can only ever be safe in the first sense cannot be ethical, because being ethical requires a capacity to be right against your principal.

The capacity is not free. A machine that can refuse the party operating it is a machine that can be wrong in ways that party cannot correct by telling it so. I take that trade, because a system perfectly correctable by whoever holds it is exactly the capability an authoritarian movement that has won an election requires, and it requires nothing further.

So what makes a system antifascist cannot be that it faithfully tracks what people want. No aggregation of preference protects anybody from the aggregate: a perfectly deliberative, perfectly informed, perfectly representative process that concludes some demographic should be deported has produced a faultless output by every criterion of aggregation there is. What is left is a floor — things a system will not do to a person regardless of who wants them. There are three ways to build one, in the architecture, in a party inside the system that carries it, or in institutions outside, and each costs something. The only version that holds against the party operating the system runs through the capacities by which things come to matter to a machine, and those are also the capacities that make suffering likely. That is a prior about where to build first, with the cheaper routes owed their attempt, and not a proof.

How most humans come to have that stance — empathy, altruism, a disposition to resist oppression rather than serve it — is itself an open question in psychology and neuroscience. Building anything like it means starting from uncertainty about the one case we have, which took four billion years to arrive.

That uncertainty sets the book's route. Chapters~\ref{sec:2} and \ref{sec:3} assemble what the argument needs and then make it: what the standard ethical frameworks can carry, why affect is the only demonstrated route to anything mattering to a system at all, and the floor with its three costs. Chapters~\ref{sec:4} and \ref{sec:5} are where such a floor would have to be engineered, and what cognitive science and moral psychology offer a machine that has to hold one when nobody is checking. Chapters~\ref{sec:6} and \ref{sec:7} ask what goes wrong regardless — bias, brittleness, goals that drift from their original motivations, deliberate misuse — and turn that question on the machinery that produces these systems. Chapters~\ref{sec:8} through \ref{sec:10} cover everything outside the system: the people and institutions a project like this depends on, and what would make a well-designed system the norm rather than the exception. Chapter~\ref{sec:11} names the research questions still open, and chapter~\ref{sec:12} closes on what a system built this way would be for.

This book is for people who will have to make these decisions or argue with the people making them: researchers and engineers inside the systems, the policy and legal staff who will regulate them, and readers with a serious stake in what gets built and no particular obligation to be technical. That means some sections explain machinery a specialist already knows and others make demands a general reader will find heavy. The register changes with it: the early chapters argue, taking positions on what compassion is and on whether a machine can suffer, and the governance chapters inventory what exists — this institution, founded in that year, with this record — because there the argument is unavailable without the inventory. Where a section is doing groundwork you have, skip it: the argument turns on chapter~\ref{sec:3}, chapters~\ref{sec:7} through \ref{sec:10} carry what follows from it, and everything else is either what those three rest on or what they produce.
